% ============================================================
% Chapter 11 — Limitations
% ============================================================
\chapter{Limitations}
\label{ch:limitations}

\roadmap{This chapter consolidates the platform's limitations as mature
scientific constraints. It groups them into measurement-validity,
data-quality, methodological, and engineering limitations, cross-referring
to where each is established. None is offered as an apology; each bounds a
claim, and several define the future work of \cref{ch:future}.}

\section{Measurement-validity limitations}
\label{sec:limitations-measurement}

\begin{enumerate}
  \item \textbf{Convergent, not criterion, validation.} RQ1's validation
  compares two machine channels and a sparse self-report; it establishes
  convergence, not agreement with ground truth
  (\cref{sec:methodology-validation}). Human-coded ground truth is
  designed-but-future (\cref{sec:sensing-coding}).
  \item \textbf{Region-level gaze.} WebGazer resolves panels, not
  fixations \citep{papoutsaki2016webgazer}; all attention claims are at
  the AOI-panel granularity (\cref{sec:sensing-gaze}).
  \item \textbf{Pupil proxy.} The pupil signal is a dark-pixel-area
  heuristic, not calibrated pupillometry, and is excluded from primary
  claims (\cref{sec:sensing-pupil}).
  \item \textbf{Placeholder SEEV weights.} The AOI \allocscore{} uses
  ordinal expected weights supplied by the researcher, not fitted SEEV
  parameters, so it is a relative alignment index rather than a validated
  SEEV measure (\cref{sec:sensing-aoi}).
  \item \textbf{No head pose.} Head-pose columns exist but are never
  populated; the report makes no head-pose claim
  (\cref{sec:sensing-facial}).
  \item \textbf{Two unwired affect derivations.} Server-windowed and
  client-per-frame affect classifiers use different weights and are not
  unified; results must name which produced them
  (\cref{sec:sensing-affect}).
  \item \textbf{Low-feasibility EF constructs.} Cognitive flexibility and
  working memory are verbal-marker proxies the implementation itself flags
  as unvalidated as executive functions
  (\cref{sec:methodology-ef}).
\end{enumerate}

\section{Data-quality limitations}
\label{sec:limitations-data}

\begin{enumerate}
  \item \textbf{Emotion CSV mislabel.} Emotion rows are labeled
  ``py-feat'' though produced by OpenFace~3.0; values are correct, the
  header is not (\cref{sec:sensing-caveats}).
  \item \textbf{Docked-scroll gap.} Sessions before the scroll-host fix
  cannot reconstruct inner reading position. \todo{Cutoff date and
  affected sessions.}
  \item \textbf{Frame truncation.} Sessions over $\sim$17 minutes recorded
  before 2026-05-21 were truncated to their earliest 5{,}000 biometric
  frames. \todo{Affected sessions.}
  \item \textbf{Silent batch drops.} Malformed byte-escape payloads cause
  whole interaction batches to be dropped with only a server warning,
  producing non-random gaps (\cref{sec:sensing-caveats}).
  \item \textbf{No CV retry.} Failed facial-analysis jobs appear as
  missing frames with no automatic recovery
  (\cref{sec:sensing-caveats}).
  \item \textbf{Sparse, late self-report.} One report per fifteen minutes,
  and none before 2026-07-01, limits the item-15 sample
  (\cref{sec:methodology-selfreport}).
  \item \textbf{Retained error rows.} EF error detections persist and must
  be filtered (\cref{sec:methodology-ef}).
\end{enumerate}

\section{Methodological limitations}
\label{sec:limitations-method}

\begin{enumerate}
  \item \textbf{No phase marker.} The database encodes no phase,
  condition, or group field; any grouped comparison rests on an externally
  defined grouping, operationalized here as pre/post
  (\cref{sec:methodology-stats}).
  \item \textbf{No stored gain or effect-size fields.} Normalized gain and
  effect sizes are computed in analysis, not stored
  (\cref{sec:methodology-gain}).
  \item \textbf{Fragile pre/post join.} Assessments are distinguished only
  by title, the attempt-to-assessment relation is unenforced, and both are
  practice-mode with multiple attempts, requiring an explicit
  attempt-selection policy (\cref{sec:methodology-gain}).
  \item \textbf{Learner-initiated confound-free but capped.} Because
  interventions are never affect-triggered, gain cannot be attributed to
  targeted support, and the design cannot test affect-contingent adaptation
  (\cref{sec:interventions-autotrigger}).
  \item \textbf{Single cohort, small $n$.} Limits power and
  generalization (\cref{sec:discussion-threats}).
\end{enumerate}

\section{Engineering limitations}
\label{sec:limitations-engineering}

\begin{enumerate}
  \item \textbf{Two unconnected KC systems.} The curated prerequisite
  graph and the mastery tracker are not wired together, blocking a
  mastery-driven adaptive path (\cref{sec:scope-item5}).
  \item \textbf{Unused BKT columns.} Mastery is an EMA; the BKT parameter
  columns are unpopulated (\cref{sec:scope-item10}).
  \item \textbf{Open AI-grading path.} Structured-question AI grading
  writes a placeholder score awaiting an external grader
  (\cref{sec:arch-maturity}).
  \item \textbf{Teacher dashboard stub and unimplemented TIMELINE studio
  output} (\cref{sec:scope-item8,sec:modes-a}).
  \item \textbf{In-memory embeddings and single-instance schedulers}
  bound scale (\cref{sec:arch-storage,sec:arch-backend}).
  \item \textbf{Export-jobs command-injection surface.} An unvalidated
  session identifier is interpolated into a shell command; the endpoint is
  role-restricted but the flaw should be fixed
  (\cref{sec:arch-maturity}).
\end{enumerate}

\medskip
\noindent Each limitation above names a specific extension. The next
chapter organizes those extensions into a program of future work.
